\chapter{CONCLUSÃO}
\label{cap:conclusao}

Este trabalho apresentou a estruturação do modelo de negócio do Grupo 3 com base em uma abordagem integrada para \textit{startups} de base tecnológica. A análise evidencia que a combinação entre estratégia adaptativa, design de plataforma, engenharia orientada a fluxo, inteligência fiscal e governança por valor fortalece a viabilidade do empreendimento desde as fases iniciais.

Como principal contribuição, o estudo organiza um caminho prático para transformar ideias em modelos de negócio testáveis e escaláveis, reduzindo incertezas típicas do processo empreendedor. Destaca-se, ainda, a importância de conectar métricas técnicas e métricas econômicas para orientar decisões de priorização com maior consistência.

Como trabalhos futuros, recomenda-se aplicar o framework a um problema real de mercado, construir o canvas completo da solução, realizar modelagem financeira com cenários e conduzir ciclos de validação com usuários e especialistas. Esses passos permitem consolidar o modelo proposto e avaliar seu potencial de implementação em contexto real.
